\section{Supervisión Legal de Pacientes y Dictámenes}

La clínica retiene legalmente el expediente fundamental. Como gerente, tienes el nivel más alto para examinar cualquier dato capturado en la base.

\subsection{Modificación de Expedientes y Doctores Externos}

Busca la sección \textbf{PatientsDashboard} en tu menú principal, luego haz clic en \textbf{Patients}.

Este menú agrupa a toda la base de datos de clientela que se ha atendido en el laboratorio. Al hacer clic sobre cualquier paciente podrás ver a fondo su ficha técnica.
\begin{itemize}
    \item \textbf{Expediente Electrónico (MRN y Raditech):} Si el área de recepción o los programas automáticos se equivocaron al generar el número de historia clínica (MRN o el Raditech ID vinculado a las facturas), puedes venir aquí a reescribirlo manualmente.
    \item \textbf{Desvincular a un Doctor Foráneo por Privacidad:} Si un paciente envía una queja o solicita explícitamente y de manera legal que uno de sus doctores tratantes pasados ya no pueda revisar ni descargar sus estudios por privacidad, baja a la parte de \textbf{Associated doctors} dentro del paciente. Aquí verás qué doctores externos tienen permiso de ver sus documentos; puedes seleccionar a cualquier doctor, borrarlo y presionar el botón "Save" (Guardar).
\end{itemize}

\subsection{Recuperación de Certificados Core}

La carpeta \textbf{Core} es donde todo culmina: aquí descansan inamovibles todas las solicitudes ("orden de compra") y los textos clínicos de estudios dictaminados.

\begin{itemize}
    \item \textbf{Verificación de Solicitudes (StudyRequests):} Si existe un problema con el dinero pagado por mostrador, entra a \textit{Study Requests}. Podrás buscar al paciente en conflicto y ver la orden inalterable con la que entró (fecha, qué estudio exacto ordenó y el doctor).
    \item \textbf{Rescate de un Reporte Textual Antiguo (Reports):} En \textit{Reports}, puedes extraer el laudo que escribió el radiólogo entero, incluyendo la línea temporal, ideal en caso de tener que acudir a soporte a instancias judiciales o de revisión.
\end{itemize}
